\documentclass{fiche}
\metadonnees{18}{Le fauteuil du dimanche}{Bloc 4 — Hypothèse, parole rapportée, traduction}

\begin{document}
\entete

\objectifs{%
\textbf{Langue} : gérondif et infinitif ; verbes de perception ; \emph{have} et \emph{get
something done}.\par
\textbf{Texte} : Katherine Mansfield, \og Miss Brill \fg (1922).\par
\textbf{Durée} : 1\,h\,30 — exercices 35 min, texte et questions 30 min, version 25 min.}

%% ===================================================================
\exercice[14 min]{Gérondif ou infinitif}

\rappel{\textbf{Rappel.} Certains verbes appellent le gérondif (\emph{-ing}), d'autres
l'infinitif ; quelques-uns acceptent les deux, mais changent alors de sens.}

\consigne{Mettez le verbe entre parenthèses à la forme qui convient.}

\begin{anglais}
\begin{enumerate}[label=\arabic*.,itemsep=2.2mm,leftmargin=*]
  \item She enjoyed \emph{(listen)} \trou{listening} to other people's conversations.
  \item She decided \emph{(wear)} \trou{to wear} her fur that afternoon.
  \item He avoided \emph{(look)} \trou{looking} at her.
  \item They stopped \emph{(talk)} \trou{to talk} to an old friend. \emph{(ils se sont
    arrêtés pour parler)}
  \item He stopped \emph{(smoke)} \trou{smoking} last year. \emph{(il a cessé)}
  \item I shall never forget \emph{(hear)} \trou{hearing} that band for the first time.
    \emph{(souvenir)}
  \item Don't forget \emph{(shut)} \trou{to shut} the box. \emph{(n'oublie pas de)}
  \item She looked forward to \emph{(hear)} \trou{hearing} them talk.
  \item It was no good \emph{(get)} \trou{getting} spectacles, she said.
  \item He suggested \emph{(buy)} \trou{buying} the gold rims.
  \item She could not help \emph{(smile)} \trou{smiling}.
  \item He refused \emph{(answer)} \trou{to answer} her.
  \item It is nice \emph{(feel)} \trou{to feel} the fur again.
  \item She went on \emph{(complain)} \trou{complaining} the whole time.
\end{enumerate}
\end{anglais}

\note[Trois listes à savoir]{Gérondif : \emph{enjoy}, \emph{avoid}, \emph{suggest},
\emph{deny}, \emph{mind}, \emph{finish}, \emph{keep}, \emph{go on}, \emph{can't help},
\emph{it's no good}, \emph{look forward to}. Infinitif : \emph{decide}, \emph{refuse},
\emph{promise}, \emph{hope}, \emph{want}, \emph{manage}, \emph{expect}, \emph{agree}. Les
deux : \emph{begin}, \emph{start}, \emph{continue}, \emph{like}, sans changement de sens.}

\note[Les quatre qui changent de sens]{\emph{Stop doing} = cesser ; \emph{stop to do} =
s'arrêter pour faire. \emph{Forget doing} = oublier un souvenir ; \emph{forget to do} =
omettre. \emph{Remember} et \emph{try} suivent le même partage : \emph{try doing} = essayer
pour voir, \emph{try to do} = s'efforcer de.}

%% ===================================================================
\exercice[10 min]{Perception et causatif}

\rappel{\textbf{Rappel.} \emph{See} et \emph{hear} + base verbale disent l'action entière,
+ \emph{-ing} l'action en cours. \emph{Have something done} : on le fait faire ; \emph{get
something done} : on s'arrange pour.}

\consigne{Complétez.}

\begin{anglais}
\begin{enumerate}[label=\arabic*.,itemsep=2.2mm,leftmargin=*]
  \item She heard the band \emph{(play)} \trou{play} the whole programme. \emph{(du début à
    la fin)}
  \item She heard the band \emph{(play)} \trou{playing} as she came down the avenue.
    \emph{(en cours)}
  \item I saw him \emph{(cross)} \trou{cross} the park and disappear.
  \item I saw him \emph{(walk)} \trou{walking} towards the rotunda.
  \item She watched the couples \emph{(pass)} \trou{pass \emph{ou} passing} in front of the
    flower-beds.
  \item She had her fur \emph{(clean)} \trou{cleaned} before the season.
  \item I must have this coat \emph{(mend)} \trou{mended}.
  \item He got his hair \emph{(cut)} \trou{cut} on Saturday.
  \item They had their photograph \emph{(take)} \trou{taken} in the gardens.
  \item She felt something \emph{(move)} \trou{move \emph{ou} moving} in her bosom.
\end{enumerate}
\end{anglais}

\note[Le causatif]{\emph{She had her fur cleaned} ne signifie pas qu'elle l'a nettoyée, mais
qu'elle l'a fait nettoyer : le participe passé marque que le sujet subit l'action au lieu de
la faire. Le français dit \og faire faire \fg, l'anglais \emph{have done}. Attention au
contresens : \emph{he had his leg broken} = on lui a cassé la jambe.}

%% ===================================================================
\exercice[11 min]{Le dimanche au jardin public}
\consigne{a) Associez chaque mot à sa traduction.}

\begin{multicols}{2}
\begin{enumerate}[label=\arabic*.,itemsep=1.2mm,leftmargin=*]
  \item a fur \hfill \trou{h}
  \item a bench \hfill \trou{c}
  \item a rotunda \hfill \trou{k}
  \item a flower-bed \hfill \trou{a}
  \item to sip \hfill \trou{f}
  \item to drift \hfill \trou{j}
  \item to stroke \hfill \trou{b}
  \item a rogue \hfill \trou{l}
  \item a tingling \hfill \trou{d}
  \item spectacles \hfill \trou{g}
  \item a rim \hfill \trou{i}
  \item to giggle \hfill \trou{e}
\end{enumerate}
\columnbreak
\begin{enumerate}[label=\alph*.,itemsep=1.2mm,leftmargin=*]
  \item un parterre de fleurs
  \item caresser
  \item un banc
  \item un picotement
  \item pouffer de rire
  \item siroter, boire à petites gorgées
  \item des lunettes
  \item une fourrure
  \item une monture (de lunettes)
  \item dériver, flotter
  \item un kiosque à musique
  \item un coquin, un garnement
\end{enumerate}
\end{multicols}

\note[Un pluriel obligatoire]{\emph{Spectacles}, \emph{glasses}, \emph{scissors},
\emph{trousers} désignent des objets faits de deux parties : toujours pluriels, et comptés
par \emph{a pair of}. Le français fait l'inverse pour \og un pantalon \fg.}

\consigne{b) Complétez avec un mot de la liste.}
\begin{anglais}
\begin{enumerate}[label=\arabic*.,itemsep=2mm,leftmargin=*]
  \item She could have taken the \trou{fur} off and \trou{stroked} it on her lap.
  \item A leaf came \trou{drifting} down from the sky.
  \item The bandsmen sat in the green \trou{rotunda} and blew out their cheeks.
  \item The Englishwoman said her \trou{spectacles} would always be sliding down her nose.
\end{enumerate}
\end{anglais}

%% ===================================================================
\newpage
\section{Texte}

\noindent\small\textit{Miss Brill, Anglaise d'un certain âge, vit seule dans une ville de
province française. Elle enseigne l'anglais et va chaque dimanche écouter la musique au
jardin public.}\normalsize

\begin{texte}
Although it was so brilliantly fine --- the blue sky powdered with gold and great spots of
light like white wine splashed over the Jardins Publiques --- Miss Brill was glad that she had
decided on her \gl[a fur]{fur}{une fourrure, ici un tour de cou}. The air was motionless, but
when you opened your mouth there was just a faint chill, like a chill from a glass of iced
water before you \gl[to sip]{sip}{boire à petites gorgées}, and now and again a leaf came
\gl[to drift]{drifting}{dériver, flotter} --- from nowhere, from the sky. Miss Brill put up her
hand and touched her fur. Dear little thing! It was nice to feel it again. She had taken it
out of its box that afternoon, shaken out the moth-powder, given it a good brush, and rubbed
the life back into the dim little eyes. ``What has been happening to me?'' said the sad little
eyes. Oh, how sweet it was to see them snap at her again from the red
\gl[an eiderdown]{eiderdown}{un édredon}! \ldots{} But the nose, which was of some black
composition, wasn't at all firm. It must have had a knock, somehow. Never mind --- a little
\gl[a dab]{dab}{une touche, un point} of black sealing-wax when the time came --- when it was
absolutely necessary\ldots{} Little \gl[a rogue]{rogue}{un coquin}! Yes, she really felt like
that about it. Little rogue biting its tail just by her left ear. She could have taken it off
and laid it on her lap and \gl[to stroke]{stroked}{caresser} it. She felt a
\gl[a tingling]{tingling}{un picotement} in her hands and arms, but that came from walking,
she supposed. And when she breathed, something light and sad --- no, not sad, exactly ---
something gentle seemed to move in her \gl[a bosom]{bosom}{la poitrine}.

There were a number of people out this afternoon, far more than last Sunday. And the band
sounded louder and gayer. That was because the Season had begun. For although the band played
all the year round on Sundays, out of season it was never the same. It was like some one
playing with only the family to listen; it didn't care how it played if there weren't any
strangers present. Wasn't the conductor wearing a new coat, too? She was sure it was new. He
scraped with his foot and flapped his arms like a \gl[a rooster]{rooster}{un coq} about to
crow, and the bandsmen sitting in the green \gl[a rotunda]{rotunda}{un kiosque à musique} blew
out their cheeks and \gl[to glare at]{glared}{fixer d'un œil furieux} at the music. Now there
came a little ``flutey'' bit --- very pretty! --- a little chain of bright drops. She was sure
it would be repeated. It was; she lifted her head and smiled.

Only two people shared her ``special'' seat: a fine old man in a velvet coat, his hands
clasped over a huge carved walking-stick, and a big old woman, sitting upright, with a roll of
knitting on her embroidered \gl[an apron]{apron}{un tablier}. They did not speak. This was
disappointing, for Miss Brill always looked forward to the conversation. She had become really
quite expert, she thought, at listening as though she didn't listen, at sitting in other
people's lives just for a minute while they talked round her.

She glanced, sideways, at the old couple. Perhaps they would go soon. Last Sunday, too, hadn't
been as interesting as usual. An Englishman and his wife, he wearing a dreadful Panama hat and
she \gl[button boots]{button boots}{des bottines à boutons}. And she'd gone on the whole time
about how she ought to wear \gl[spectacles]{spectacles}{des lunettes}; she knew she needed
them; but that it was no good getting any; they'd be sure to break and they'd never keep on.
And he'd been so patient. He'd suggested everything --- gold \gl[a rim]{rims}{une monture},
the kind that curved round your ears, little pads inside the bridge. No, nothing would please
her. ``They'll always be sliding down my nose!'' Miss Brill had wanted to shake her.
\end{texte}
\source{Katherine Mansfield, \og Miss Brill \fg, \emph{The Garden Party and Other Stories},
Londres, Constable, 1922.}

\partiel[15 min]{Questions}
\consigne{\en{Answer in English, in full sentences.}}

\begin{enumerate}[label=\arabic*.,itemsep=2mm,leftmargin=*]
  \item \en{How does Miss Brill treat her fur?}
    \lignes{3}
    \corr{Like a small living creature and a friend: she brushes it, rubs the life back into
    its eyes, lends it a voice, calls it \emph{dear little thing} and \emph{little rogue},
    and would like to stroke it on her lap.}
  \item \en{What does she notice about the band, and how does she explain it?}
    \lignes{3}
    \corr{That it sounds louder and gayer, and that the conductor is wearing a new coat. She
    explains it by the Season having begun: out of season the band plays as if only the
    family were listening.}
  \item \en{What is her ``special'' seat for?}
    \lignes{3}
    \corr{For listening. She sits there every Sunday to overhear the conversations of the
    people who share the bench, and has become expert at \emph{listening as though she
    didn't listen}.}
  \item \en{Why is she disappointed by the old couple, and what was wrong with last Sunday?}
    \lignes{3}
    \corr{The old couple do not speak at all. The Sunday before, the English couple talked,
    but only about spectacles, and the wife refused every solution her patient husband
    offered.}
  \item \en{``Miss Brill had wanted to shake her.'' What does this reaction reveal?}
    \lignes{3}
    \corr{That she takes the lives she overhears personally, as if she had a part in them.
    She is not a spectator but a participant who has not been told she is one.}
  \item \en{Find three details that suggest Miss Brill is lonely without ever saying so.}
    \lignes{3}
    \corr{The fur she talks to; the fur kept in a box with moth-powder, taken out for the
    occasion; the \emph{special} seat; her expertise at sitting in other people's lives
    \emph{just for a minute}; the fact that nobody in the whole scene speaks to her.}
\end{enumerate}

\note[Gérondifs et infinitifs du texte]{\emph{glad that she had decided on her fur},
\emph{It was nice to feel it}, \emph{She could have taken it off}, \emph{looked forward to the
conversation}, \emph{expert at listening}, \emph{at sitting in other people's lives},
\emph{it was no good getting any}, \emph{He'd suggested everything}, \emph{nothing would
please her}. Après une préposition (\emph{at}, \emph{forward to}, \emph{no good}), toujours
le gérondif : c'est la règle sans exception.}

%% ===================================================================
\newpage
\section{Version}
\consigne{Traduisez la fin de la nouvelle. Entre-temps, Miss Brill a compris que tous ces
gens, elle comprise, jouaient une pièce, et elle en a eu les larmes aux yeux.}

\begin{version}
Just at that moment a boy and girl came and sat down where the old couple had been. They were
beautifully dressed; they were in love. The hero and heroine, of course, just arrived from his
father's yacht. And still soundlessly singing, still with that trembling smile, Miss Brill
prepared to listen.

``No, not now,'' said the girl. ``Not here, I can't.''

``But why? Because of that stupid old thing at the end there?'' asked the boy. ``Why does she
come here at all --- who wants her? Why doesn't she keep her silly old
\gl[a mug]{mug}{une figure, une tête (familier)} at home?''

``It's her fu-fur which is so funny,'' \gl[to giggle]{giggled}{pouffer} the girl. ``It's
exactly like a fried \gl[a whiting]{whiting}{un merlan}.''

``Ah, be off with you!'' said the boy in an angry whisper. Then: ``Tell me, ma petite
chère ---''

``No, not here,'' said the girl. ``Not \emph{yet}.''

On her way home she usually bought a slice of honey-cake at the baker's. It was her Sunday
treat. Sometimes there was an almond in her slice, sometimes not. It made a great difference.
If there was an almond it was like carrying home a tiny present --- a surprise --- something
that might very well not have been there. She hurried on the almond Sundays and struck the
match for the kettle in quite a \gl[dashing]{dashing}{fringant, allègre} way.

But to-day she passed the baker's by, climbed the stairs, went into the little dark room ---
her room like a cupboard --- and sat down on the red eiderdown. She sat there for a long time.
The box that the fur came out of was on the bed. She unclasped the
\gl[a necklet]{necklet}{un tour de cou} quickly; quickly, without looking, laid it inside. But
when she put the lid on she thought she heard something crying.
\end{version}
\source{\emph{Ibid.}}

\lignes{22}

\corr{\textbf{Proposition de traduction.} À cet instant précis, un garçon et une fille vinrent
s'asseoir à la place du vieux couple. Ils étaient magnifiquement habillés ; ils étaient
amoureux. Le héros et l'héroïne, naturellement, débarquant du yacht de son père. Et, chantant
toujours sans un son, toujours avec ce sourire tremblant, Miss Brill s'apprêta à écouter.

\og Non, pas maintenant, dit la fille. Pas ici, je ne peux pas. \fg

\og Mais pourquoi ? À cause de cette vieille bêtise là-bas, au bout ? demanda le garçon.
Pourquoi vient-elle seulement ici --- qui veut d'elle ? Pourquoi ne garde-t-elle pas sa vieille
figure ridicule à la maison ? \fg

\og C'est sa fou-fourrure qui est si drôle, pouffa la fille. On dirait exactement un merlan
frit. \fg

\og Allons, laisse donc ! \fg{} dit le garçon dans un murmure agacé. Puis : \og Dis-moi, ma
petite chère\ldots \fg

\og Non, pas ici, dit la fille. Pas \emph{encore}. \fg

En rentrant, elle achetait d'ordinaire une tranche de gâteau au miel chez le boulanger.
C'était sa gâterie du dimanche. Il y avait parfois une amande dans sa tranche, parfois non.
Cela faisait une grande différence. S'il y avait une amande, c'était comme de rapporter chez
soi un tout petit cadeau --- une surprise --- quelque chose qui aurait fort bien pu ne pas s'y
trouver. Les dimanches à l'amande, elle pressait le pas et craquait l'allumette sous la
bouilloire d'un air tout à fait allègre.

Mais aujourd'hui elle passa devant la boulangerie sans s'arrêter, monta l'escalier, entra dans
la petite pièce sombre --- sa chambre pareille à un placard --- et s'assit sur l'édredon rouge.
Elle resta là longtemps. La boîte d'où la fourrure était sortie était sur le lit. Elle
dégrafa vivement le tour de cou ; vivement, sans regarder, le posa à l'intérieur. Mais quand
elle remit le couvercle, il lui sembla entendre quelque chose qui pleurait.}

\note[Choix de traduction 1]{\emph{something that might very well not have been there} :
\emph{might have} + participe à la forme négative (fiche 10) : \og qui aurait fort bien pu ne
pas s'y trouver \fg. Ordre des mots à surveiller : la négation porte sur l'infinitif, pas
sur le modal.}

\note[Choix de traduction 2]{\emph{she passed the baker's by} : verbe à particule séparé par
son complément. \emph{To pass by} = passer devant sans s'arrêter --- et c'est tout le
sens de la fin : elle renonce à sa gâterie.}

\note[Choix de traduction 3]{\emph{the baker's} : génitif elliptique, le nom du commerce est
sous-entendu (\emph{shop}). Comme \emph{at the butcher's}, \emph{to the doctor's}. Le français
dit \og chez le boulanger \fg.}

\note[Choix de traduction 4]{\emph{she thought she heard something crying} : \emph{hear} +
complément + \emph{-ing}, action en cours (exercice 2). \og Il lui sembla entendre quelque
chose qui pleurait \fg{} : le français passe par une relative. Et le dernier mot de la
nouvelle est \emph{something}, non \emph{someone} : c'est la fourrure qui pleure, ou bien
ce n'est pas la fourrure.}

%% ===================================================================
\partiel{Lexique à mémoriser}
\begin{itemize}[label=--,itemsep=0.8mm,leftmargin=*]\small
  \item \textbf{a fur} — une fourrure ; \textit{adjectif} furry
  \item \textbf{a bench} — un banc
  \item \textbf{a lap} — les genoux (assis) ; \textit{à distinguer de} a knee
  \item \textbf{to stroke} — caresser ; \textit{et} un coup, une attaque
  \item \textbf{to sip} — boire à petites gorgées
  \item \textbf{a chill} — un froid léger \textit{(fiche 10)}
  \item \textbf{to drift} — dériver ; \textit{nom} a drift
  \item \textbf{a rogue} — un coquin
  \item \textbf{to snap} — claquer ; \textit{ici} briller vivement
  \item \textbf{a knock} — un coup, un choc ; \textit{verbe} to knock
  \item \textbf{to flap} — battre (des ailes, des bras)
  \item \textbf{to scrape} — racler ; \textit{et} gratter
  \item \textbf{to glance at} — jeter un coup d'œil à ; \textit{à distinguer de} to gaze \textit{(fiche 16)}
  \item \textbf{sideways} — de côté ; side + ways\textit{, comme} always, otherwise
  \item \textbf{to look forward to} — attendre avec impatience \textit{(+ gérondif)}
  \item \textbf{to go on about} — rabâcher, ne pas cesser de parler de
  \item \textbf{it is no good} — cela ne sert à rien \textit{(+ gérondif)}
  \item \textbf{a treat} — une gâterie, un plaisir ; \textit{verbe} to treat
  \item \textbf{an almond} — une amande ; \textit{le} l \textit{ne se prononce pas}
  \item \textbf{a lid} — un couvercle ; an eyelid \textit{(fiche 15)}
\end{itemize}

\end{document}
